\documentclass[12pt,a4paper]{article}
\usepackage[spanish]{babel}
\usepackage[utf8]{inputenc}
\usepackage[T1]{fontenc}
\usepackage{lmodern}
\usepackage[table]{xcolor}
\usepackage{geometry}
\usepackage{setspace}
\usepackage{longtable}
\usepackage{booktabs}
\usepackage{array}
\usepackage{tabularx}
\usepackage{hyperref}
\usepackage{enumitem}

\geometry{margin=2.4cm}
\setstretch{1.12}
\setlength{\parindent}{0pt}
\setlength{\parskip}{0.55em}
\hypersetup{colorlinks=true, linkcolor=black, urlcolor=blue, citecolor=black}

\definecolor{PrimaryBlue}{HTML}{16324F}
\definecolor{SoftGray}{HTML}{F3F6F8}

	itle{\Huge\textbf{Casos de Uso}\[0.3em]\Large Plataforma Web de Mascotas 3D}
\author{Bryan Quispe}
\date{10 de mayo de 2026}

\begin{document}

\begin{titlepage}
\centering
\vspace*{1.8cm}
{\Huge\bfseries\color{PrimaryBlue} Casos de Uso\\[0.35em]}
{\LARGE Especificaci\'on detallada de flujos de negocio\\[1.2cm]}

\begin{tabular}{p{5cm} p{8cm}}
\rowcolor{SoftGray}
\textbf{Versi\'on} & 1.0 \\
\textbf{Fecha} & 10 de mayo de 2026 \\
\textbf{Autor} & Bryan Quispe \\
\textbf{Alcance} & MVP en desarrollo: Flujos principales \\
\end{tabular}

\vfill
{\large Documento para compilaci\'on en Overleaf}
\end{titlepage}

\tableofcontents
\newpage

\section{Introducci\'on}
Este documento describe los casos de uso principales de la plataforma en desarrollo. Un caso de uso representa la interacci\'on entre un actor humano y el sistema dentro de una frontera funcional concreta. Cada caso incluye acciones del usuario, respuestas del sistema, precondiciones, flujos alternos y postcondiciones.

\section{CU-01: Registrarse en la plataforma}

\begin{tabular}{|p{3.5cm}|p{10cm}|}
\hline
\textbf{Identificador} & CU-01 \\
\hline
\textbf{Nombre} & Registrarse en la plataforma \\
\hline
\textbf{Descripci\'on} & Un nuevo usuario se crea una cuenta en el sistema. \\
\hline
\textbf{Actor principal} & Usuario no autenticado \\
\hline
\textbf{Precondiciones} & El usuario accede al formulario de registro. \\
\hline
\end{tabular}

\vspace{0.5em}

\textbf{Flujo normal}:
\begin{enumerate}[leftmargin=*]
    \item El usuario ingresa email, contrase\~na y confirmaci\'on de contrase\~na.
    \item El sistema valida que el email sea \'unico.
    \item El sistema valida que la contrase\~na cumpla requisitos (8+ caracteres).
    \item El sistema crea el usuario en la base de datos con contrase\~na hasheada.
    \item El sistema retorna mensaje de \'exito.
    \item El usuario es redirigido a la p\'agina de login.
\end{enumerate}

\textbf{Flujos alternativos}:
\begin{itemize}[leftmargin=*]
    \item FA-01: Email ya existe $\rightarrow$ Mostrar error, permitir reintento.
    \item FA-02: Contrase\~na d\'ebil $\rightarrow$ Mostrar requisitos, permitir reintento.
    \item FA-03: Contrase\~nas no coinciden $\rightarrow$ Mostrar error, permitir reintento.
\end{itemize}

\textbf{Postcondiciones}:
\begin{itemize}[leftmargin=*]
    \item Usuario registrado en la base de datos.
    \item Usuario puede iniciar sesi\'on.
\end{itemize}

\begin{figure}[h]
\centering
\fbox{\parbox[c][3cm][c]{0.6\textwidth}{\centering Diagrama de caso de uso - CU-01}}
\caption{Diagrama conceptual: CU-01 Registrarse}
\end{figure}

\newpage

\section{CU-02: Iniciar sesi\'on}

\begin{tabular}{|p{3.5cm}|p{10cm}|}
\hline
\textbf{Identificador} & CU-02 \\
\hline
\textbf{Nombre} & Iniciar sesi\'on \\
\hline
\textbf{Descripci\'on} & Un usuario registrado inicia sesi\'on y obtiene acceso a su cuenta. \\
\hline
\textbf{Actor principal} & Usuario no autenticado \\
\hline
\textbf{Precondiciones} & El usuario est\'a registrado y accede al formulario de login. \\
\hline
\end{tabular}

\vspace{0.5em}

\textbf{Flujo normal}:
\begin{enumerate}[leftmargin=*]
    \item El usuario ingresa email y contrase\~na.
    \item El sistema valida credenciales contra la base de datos.
    \item El sistema genera un JWT v\'alido por 24 horas.
    \item El sistema retorna el token al frontend.
    \item El frontend guarda el token en localStorage.
    \item El usuario es redirigido al dashboard.
\end{enumerate}

\textbf{Flujos alternativos}:
\begin{itemize}[leftmargin=*]
    \item FA-01: Credenciales inv\'alidas $\rightarrow$ Mostrar error gen\'erico, permitir reintento.
    \item FA-02: Usuario no existe $\rightarrow$ Mostrar error gen\'erico, permitir reintento.
\end{itemize}

\textbf{Postcondiciones}:
\begin{itemize}[leftmargin=*]
    \item Usuario autenticado con sesi\'on activa.
    \item Token almacenado en el cliente.
\end{itemize}

\begin{figure}[h]
\centering
\fbox{\parbox[c][3cm][c]{0.6\textwidth}{\centering Diagrama de caso de uso - CU-02}}
\caption{Diagrama conceptual: CU-02 Login}
\end{figure}

\newpage

\section{CU-03: Registrar una mascota}

\begin{tabular}{|p{3.5cm}|p{10cm}|}
\hline
\textbf{Identificador} & CU-03 \\
\hline
\textbf{Nombre} & Registrar una mascota \\
\hline
\textbf{Descripci\'on} & Un usuario autenticado crea un registro de mascota. \\
\hline
\textbf{Actor principal} & Usuario autenticado \\
\hline
\textbf{Precondiciones} & El usuario ha iniciado sesi\'on y accede a la p\'agina de crear mascota. \\
\hline
\end{tabular}

\vspace{0.5em}

\textbf{Flujo normal}:
\begin{enumerate}[leftmargin=*]
    \item El usuario ingresa nombre, edad y tipo de animal.
    \item El sistema valida campos obligatorios.
    \item El sistema crea el registro de mascota asociado al usuario.
    \item El sistema retorna ID de la mascota.
    \item El usuario ve confirmaci\'on de \'exito.
    \item El usuario es redirigido a la lista de mascotas.
\end{enumerate}

\textbf{Flujos alternativos}:
\begin{itemize}[leftmargin=*]
    \item FA-01: Campos obligatorios incompletos $\rightarrow$ Mostrar validaci\'on, permitir correcci\'on.
    \item FA-02: Token expirado $\rightarrow$ Redirigir a login.
\end{itemize}

\textbf{Postcondiciones}:
\begin{itemize}[leftmargin=*]
    \item Mascota registrada en la base de datos.
    \item Asociada al usuario propietario.
    \item Visible en el listado del usuario.
\end{itemize}

\begin{figure}[h]
\centering
\fbox{\parbox[c][3cm][c]{0.6\textwidth}{\centering Diagrama de caso de uso - CU-03}}
\caption{Diagrama conceptual: CU-03 Registrar mascota}
\end{figure}

\newpage

\section{CU-04: Ver lista de mascotas}

\begin{tabular}{|p{3.5cm}|p{10cm}|}
\hline
\textbf{Identificador} & CU-04 \\
\hline
\textbf{Nombre} & Ver lista de mascotas \\
\hline
\textbf{Descripci\'on} & Usuario autenticado visualiza todas sus mascotas. \\
\hline
\textbf{Actor principal} & Usuario autenticado \\
\hline
\textbf{Precondiciones} & El usuario ha iniciado sesi\'on. \\
\hline
\end{tabular}

\vspace{0.5em}

\textbf{Flujo normal}:
\begin{enumerate}[leftmargin=*]
    \item El usuario accede al dashboard.
    \item El sistema consulta mascotas del usuario autenticado.
    \item El sistema retorna lista en JSON con id, nombre, edad, tipo.
    \item El frontend renderiza tarjetas con informaci\'on de cada mascota.
    \item El usuario puede clickear en una mascota para ver detalles.
\end{enumerate}

\textbf{Flujos alternativos}:
\begin{itemize}[leftmargin=*]
    \item FA-01: No hay mascotas registradas $\rightarrow$ Mostrar mensaje vac\'io y bot\'on para crear.
    \item FA-02: Token expirado $\rightarrow$ Redirigir a login.
\end{itemize}

\textbf{Postcondiciones}:
\begin{itemize}[leftmargin=*]
    \item Usuario ve su lista de mascotas.
    \item Puede navegar a detalles de cada mascota.
\end{itemize}

\begin{figure}[h]
\centering
\fbox{\parbox[c][3cm][c]{0.6\textwidth}{\centering Diagrama de caso de uso - CU-04}}
\caption{Diagrama conceptual: CU-04 Listado de mascotas}
\end{figure}

\newpage

\section{CU-05: Ver detalle de mascota y modelo 3D}

\begin{tabular}{|p{3.5cm}|p{10cm}|}
\hline
\textbf{Identificador} & CU-05 \\
\hline
\textbf{Nombre} & Ver detalle de mascota y modelo 3D \\
\hline
\textbf{Descripci\'on} & Usuario autenticado visualiza el detalle de una mascota y su modelo 3D dentro del sistema. \\
\hline
\textbf{Actor principal} & Usuario autenticado \\
\hline
\textbf{Precondiciones} & El usuario ha iniciado sesi\'on y selecciona una mascota del listado. \\
\hline
\end{tabular}

\vspace{0.5em}

\textbf{Flujo normal}:
\begin{enumerate}[leftmargin=*]
    \item El usuario selecciona una mascota espec\'ifica desde la lista.
    \item El sistema valida que la mascota pertenezca al usuario autenticado.
    \item El sistema recupera los datos de la mascota, su foto y el modelo 3D asociado.
    \item El sistema renderiza el detalle en pantalla.
    \item El usuario interact\'ua con el modelo para verlo desde distintos \u00e1ngulos.
    \item El usuario puede regresar al listado o continuar con otra mascota.
\end{enumerate}

\textbf{Flujos alternativos}:
\begin{itemize}[leftmargin=*]
    \item FA-01: Modelo 3D no disponible $\rightarrow$ Mostrar placeholder.
    \item FA-02: Foto no existe $\rightarrow$ Mostrar imagen por defecto.
    \item FA-03: Token expirado $\rightarrow$ Redirigir a login.
\end{itemize}

\textbf{Postcondiciones}:
\begin{itemize}[leftmargin=*]
    \item Usuario visualiza detalle de su mascota.
    \item Modelo 3D renderizado correctamente.
\end{itemize}

\begin{figure}[h]
\centering
\fbox{\parbox[c][3cm][c]{0.6\textwidth}{\centering Diagrama de caso de uso - CU-05}}
\caption{Diagrama conceptual: CU-05 Detalle y modelo 3D}
\end{figure}

\newpage

\section{CU-06: Expiraci\'on de sesi\'on por inactividad}

\begin{tabular}{|p{3.5cm}|p{10cm}|}
\hline
\textbf{Identificador} & CU-06 \\
\hline
\textbf{Nombre} & Expiraci\'on de sesi\'on por inactividad \\
\hline
\textbf{Descripci\'on} & El sistema cierra la sesi\'on tras inactividad prolongada. \\
\hline
\textbf{Actor principal} & Usuario autenticado \\
\hline
\textbf{Precondiciones} & El usuario est\'a autenticado y mantiene una sesi\'on activa. \\
\hline
\end{tabular}

\vspace{0.5em}

\textbf{Flujo normal}:
\begin{enumerate}[leftmargin=*]
    \item El usuario interact\'ua con el sistema y reinicia el contador de inactividad.
    \item El sistema detecta la falta de actividad durante el umbral definido.
    \item El sistema invalida la sesi\'on y elimina credenciales locales.
    \item El usuario es redirigido al inicio de sesi\'on al intentar continuar.
\end{enumerate}

\textbf{Postcondiciones}:
\begin{itemize}[leftmargin=*]
    \item Usuario no autenticado.
    \item Token removido del cliente.
\end{itemize}

\begin{figure}[h]
\centering
\fbox{\parbox[c][3cm][c]{0.6\textwidth}{\centering Diagrama de caso de uso - CU-06}}
\caption{Diagrama conceptual: CU-06 Expiraci\'on de sesi\'on por inactividad}
\end{figure}

\newpage
\section{Referencias}
Fuentes consultadas para formato y convenciones de casos de uso:
\begin{thebibliography}{9}
\bibitem{scrumguide}
K. Schwaber y J. Sutherland,
\textit{The Scrum Guide: The Definitive Guide to Scrum: The Rules of the Game},
The Scrum Guide, 2020. Dispon\'ible en: \url{https://scrumguides.org}

\bibitem{ieee830}
IEEE,
\textit{IEEE Recommended Practice for Software Requirements Specifications},
IEEE Std 830-1998.

\bibitem{cockburn}
Alistair Cockburn,
\textit{Writing Effective Use Cases}, Addison-Wesley, 2000.
\end{thebibliography}

\end{document}
